\par\Needspace{6\baselineskip}\medskip
\hypertarget{native-124-article091}{}\pdfbookmark[2]{91.}{article-091}
\begin{center}91.\end{center}
In the case under consideration, the upper sign in equation 12 must clearly be taken. Introducing the quantity $\xi$ puts this equation in the form
\[
m=(l+x)^{\frac12}+\dfrac{(l+x)^{\frac32}}{\frac34-\frac{9}{10}(x-\xi)}\mathpunct{\cdot}
\]
Thus, by setting $\BareRoot{(l+x)}=\dfrac{m}{y}\mathpunct{\raisebox{.45ex}{,}}$ and
\[
\mathrm{14)}\qquad \dfrac{mm}{\frac56+l+\xi}=h,
\]
and carrying out the reductions, we obtain
\[
\mathrm{15)}\qquad h=\dfrac{(y-1)yy}{y+\frac13}\mbox{\HistoricalDisplayNoteMark}.
\]
\noindent\HistoricalDisplayNoteText{Manuscript note:]
\[
y-1=\dfrac{l+x}{\frac34-\frac{9}{10}(x-\xi)}
\]
\[
y=1+\frac{10}{9}h-\frac{110}{81}hh+\frac{2410}{729}h^3-\text{etc.}
\]
\[
yy=1+\frac{20}{9}h-\frac{40}{27}hh+\frac{2620}{729}h^3-\text{etc.}
\]
\[
\dfrac{1}{yy}=1-\frac{20}{9}h+\frac{520}{81}hh-\frac{5140}{243}h^3+\text{etc.}
\]
}If we may therefore regard $h$ as known, the cubic equation determines $y$, and then we have
\[
\mathrm{16)}\qquad x=\dfrac{mm}{yy}-l.
\]
Although $h$ involves the still unknown quantity $\xi$, we may neglect it in the \SourceLeaf{125}first approximation and take $h$ to be $\dfrac{mm}{\frac56+l}\mathpunct{\raisebox{.45ex}{,}}$ since $\xi$ is certainly always very small in the case under consideration. Equations 15 and 16 then give $y$ and $x$. From $x$, Table III gives $\xi$, with which formula 14 gives the corrected value of $h$. Repeating the same calculation with this value gives corrected values of $y$ and $x$. These generally differ so little from their previous values that $\xi$, taken again from Table III, has the same value as before. Otherwise the calculation must be repeated until this value ceases to change. As soon as $x$ has been found, $g$ is obtained from the formula $\sin\tfrac12 g^2=x$.

These rules apply to the first case, where $\cos f$ is positive. In the other case, where it is negative, we set $\BareRoot{(L-x)}=\dfrac{M}{Y}$ and
\[
\mathrm{14}^{\ast}\mathrm{)}\qquad \dfrac{MM}{L-\frac56-\xi}=H,
\]
so that equation $12^{\ast}$, after reduction, becomes
\[
\mathrm{15}^{\ast}\mathrm{)}\qquad H=\dfrac{(Y+1)YY}{Y-\frac13}\mathpunct{\cdot}
\]
This cubic equation thus determines $Y$ from $H$, and $x$ is in turn obtained from the equation
\[
\mathrm{16}^{\ast}\mathrm{)}\qquad x=L-\dfrac{MM}{YY}\mathpunct{\cdot}
\]
For the first approximation we take $H$ to have the value $\dfrac{MM}{L-\frac56}\mathpunct{\raisebox{.45ex}{;}}$ using the value of $x$ obtained from it by equations $15^{\ast}$ and $16^{\ast}$, we take $\xi$ from Table III. Formula $14^{\ast}$ then gives the corrected value of $H$, with which we repeat the calculation in the same way. Finally, from $x$ we determine the angle $g$ in the same way as in the first case.
